\section{The direct iPRA bi-interpretation}\label{sec:biinterp}
\subsection{The equation-induction to iPRA step}

At this point the logical interface supplies ordinary intuitionistic logic,
coverage $\forall x(x=0\lor\exists y\,x=Sy)$, separation, and induction
for literal equations with arbitrary retained side assumptions. These are
exactly the assumptions used in the following internal arithmetic argument.

Using $\operatorname{Rec}$, define addition by $x+0=x$ and
$x+Sy=S(x+y)$, and truncated subtraction by
$d(x,0)=x$, $d(x,Sy)=p(d(x,y))$. Equation induction proves
$0+y=y$ and $Sx+y=S(x+y)$. Equation induction in $y$ in
\[
 p(d(Sx,y))=d(x,y)
\]
has base $p(Sx)=x$ and step obtained by applying $p$ to its hypothesis.
It gives $d(Sx,Sy)=d(x,y)$, hence a further equation induction gives
$d(Sa,a)=1$.

The remainder equation
\[
 C(\operatorname{sg}(d(x,y)),x,y+d(x,y))=x
\tag{Rem}
\]
also follows by equation induction on $y$, keeping $x$ fixed. At zero,
both selector alternatives equal $x$. At the step put $r=d(x,y)$ and
$r'=p(r)$. Coverage splits $r'=0$ from $r'=Sz$. In the zero case
the new equation selects $x=x$. In the successor case, coverage of $r$,
the predecessor equations, and separation imply $r=Sr'$. The old equation
then gives $y+r=x$, and the new selected sum is
\[
 Sy+r'=S(y+r')=y+Sr'=y+r=x.
\]
The branch distinctions use coverage, not induction on a disjunction.
The only induction hypothesis throughout is the single equation (Rem).

If $d(Sa,b)=1$, (Rem) gives $b+1=Sa$ and hence $a=b$ by
successor injectivity. The reverse implication uses $d(Sa,a)=1$.
Put
\[
 K(n)=C(n,1,p(n)),\qquad
 \epsilon(a,b)=\operatorname{sg}(K(d(Sa,b))).
\]
Coverage and separation prove $K(n)=0\leftrightarrow n=1$ and
$\operatorname{sg}(n)=0\leftrightarrow n=0$. Therefore
\[
 \epsilon(a,b)=0\leftrightarrow a=b,
 \qquad \epsilon(a,b)=0\lor\epsilon(a,b)=1.
\]
This is a whole-domain theorem, not an experiment on standard numerals.

Every fixed quantifier-free formula has a Boolean characteristic term:
\[
\begin{aligned}
 b_{s=t}&=\epsilon(s,t),& b_\bot&=1,\\
 b_{A\land B}&=C(b_A,b_B,1),&
 b_{A\lor B}&=C(b_A,0,b_B),&
 b_{A\to B}&=C(b_A,b_B,0).
\end{aligned}
\]
External induction on formula syntax, using the two Boolean cases, proves
$b_A=0\leftrightarrow A$. In the implication case only the derived
decidability of the quantifier-free antecedent is used. Equation induction
on $b_{A(n)}=0$ gives precisely the quantifier-free induction schema of
iPRA, with the same side assumptions.

\subsection{The two whole-domain interpretations}

Let $T_R$ be ordinary intuitionistic first-order logic with equality,
the recursor computations, coverage, separation, and contextual equation
induction. Interpret each fixed body $h(i,u,\vec p)$ in iPRA by its
primitive-recursive description, and interpret its recursor by
\[
 r_h(0,a,\vec p)=Sa,\qquad
 r_h(Sn,a,\vec p)=h(n,r_h(n,a,\vec p),\vec p).
\]
This gives $\alpha:T_R\to\mathrm{iPRA}$. Canonical zero nodes are
translated by successor. Coverage in iPRA follows by quantifier-free
induction on $x=0\lor x=S(p(x))$, where $p$ is its standard
primitive-recursive predecessor: zero takes the left disjunct and each
successor takes the right disjunct by $p(Sn)=n$. Existential introduction
then gives the stated coverage axiom. All other axioms translate immediately.

Conversely compile each fixed primitive-recursive description by the
displayed $\operatorname{Rec}$ to obtain
$\beta:\mathrm{iPRA}\to T_R$. Its defining equations, successor axioms,
and complete quantifier-free induction have all been proved above.
Both maps leave the numerical domain and actual equality unchanged.

Pointwise equal fixed update bodies and equal seeds give equal recursors:
induct on a fresh counter in this one equation. The zero case uses
successor congruence on the seeds. The successor case replaces the preceding
state by the induction equation and then uses the pointwise update equality.
This body-extensionality lemma handles different syntactic descriptions of
extensionally equal update functions; they need not be identified as symbols.

For an $R$ term, external induction on its finite syntax identifies its
children and update body with their translated composites. One equation
induction compares the original and compiled trajectories, with identical
seed $Sa$ and pointwise equal updates. Thus $T_R\vdash t^{\alpha\beta}=t$.
For a primitive-recursive description, external induction on that description
does the converse. A recursion case again uses one equation induction to
identify two trajectories with equal bases and steps. Hence
$\mathrm{iPRA}\vdash t^{\beta\alpha}=t$.
Induction on ordinary formula syntax lifts these equalities to both formula
round trips. The internally definable numerical composite isomorphisms have
graph $x=y$.

\subsection{Normalization does not hide an arithmetic assumption}\label{sec:normalization-soundness}

Canonical normalization recursively normalizes all children, and replaces
$R(0,a;i,u.h)$ by $R(0,a;i,u.0)$. It is a terminating finite-tree
traversal; it performs no successor-counter unfolding. It is idempotent,
scope preserving, and obeys
\[
 N(t\sigma)=N(N(t)N(\sigma)).
\]
Once a counter has become literal zero, subsequent substitution cannot make
its erased body relevant again. This gives the substitution identity by
structural induction. Under binders the usual lifting conventions are used.

In the raw ordinary theory, $N(t)=t$ is proved by external term induction.
Use equality congruence for the counter and initializer, the preceding
body-extensionality lemma for a changed update body, and the raw zero
equation at an erased body. Thus normalization is justified even beneath
arbitrary enclosing recursors. The argument uses already available equation
induction in the ordinary theory, not an unproved equality principle in the
new checker. Its proof-theoretic converse is supplied by the recovered
computation, equality, and equation-induction rules of the canonical calculus.

These are written proofs. The existing Lean formalization of the scalar
calculus does not by itself certify this different counter-aware kernel,
its canonical substitution code, or a new checker implementation.
